\documentclass[11pt,a4paper]{scrartcl}
% =====================================================================
%  pds-apuntes.tex — Estilo común de apuntes y hojas de ejercicios.
%
%  Uso (desde teoría/Tx/apuntes/ o teoría/Tx/ejercicios/):
%     \documentclass[11pt,a4paper]{scrartcl}   % o scrreprt
%     % =====================================================================
%  pds-apuntes.tex — Estilo común de apuntes y hojas de ejercicios.
%
%  Uso (desde teoría/Tx/apuntes/ o teoría/Tx/ejercicios/):
%     \documentclass[11pt,a4paper]{scrartcl}   % o scrreprt
%     % =====================================================================
%  pds-apuntes.tex — Estilo común de apuntes y hojas de ejercicios.
%
%  Uso (desde teoría/Tx/apuntes/ o teoría/Tx/ejercicios/):
%     \documentclass[11pt,a4paper]{scrartcl}   % o scrreprt
%     \input{../../comun/pds-apuntes}
%     \tema{1}{Introducción: señales y sistemas}
%     \begin{document} \cabecera{Apuntes} ... \end{document}
%  Compilar con:  latexmk -xelatex  (o «make» desde teoría/)
% =====================================================================

\input{../../comun/pds-comun.tex}

\usepackage[a4paper, top=2.6cm, bottom=2.6cm, left=2.5cm, right=2.5cm,
            headheight=1.1cm]{geometry}
\usepackage{fontspec}
\usepackage{unicode-math}
\setmainfont{LibertinusSerif}[Extension=.otf, UprightFont=*-Regular, BoldFont=*-Semibold,
  ItalicFont=*-Italic, BoldItalicFont=*-SemiboldItalic]
\setsansfont{FiraSans}[Extension=.otf, UprightFont=*-Regular, BoldFont=*-SemiBold,
  ItalicFont=*-Italic, BoldItalicFont=*-SemiBoldItalic, Scale=MatchLowercase]
\setmonofont{FiraMono}[Extension=.otf, UprightFont=*-Regular, BoldFont=*-Medium, AutoFakeSlant=0.15, Scale=0.85]
\setmathfont{LibertinusMath-Regular.otf}

\usepackage{microtype}
\usepackage{enumitem}
\usepackage{multicol}
\setlist{itemsep=2pt, topsep=4pt}
\usepackage{caption}
\captionsetup{font=small, labelfont={sf,bf,color=pdsazul}, labelsep=period}
\usepackage{subcaption}
\usepackage[most]{tcolorbox}
\usepackage{environ}
\usepackage{scrlayer-scrpage}
\usepackage[hidelinks, unicode]{hyperref}
\usepackage[spanish, noabbrev]{cleveref}

% Títulos en Fira Sans con el azul de la asignatura
\setkomafont{disposition}{\sffamily\bfseries\color{pdsazul}}
\setkomafont{descriptionlabel}{\sffamily\bfseries}
\setlength{\parindent}{0pt}
\setlength{\parskip}{0.55\baselineskip plus 2pt}
\linespread{1.05}

% ---------------------------------------------------------------------
%  Datos del tema y cabeceras
% ---------------------------------------------------------------------
\newcommand{\numtema}{}
\newcommand{\nombretema}{}
\newcommand{\tema}[2]{\renewcommand{\numtema}{#1}\renewcommand{\nombretema}{#2}}

\pagestyle{scrheadings}
\clearpairofpagestyles
\ihead{\sffamily\small\color{pdsgris}\asignatura}
\ohead{\sffamily\small\color{pdsgris}Tema~\numtema. \nombretema}
\ifoot{\sffamily\small\color{pdsgris}\grado}
\ofoot{\sffamily\small\pagemark}
\setkomafont{pagehead}{\normalfont}
\setkomafont{pagefoot}{\normalfont}

% \cabecera{Apuntes|Ejercicios|...} : bloque de título de la primera página
\newcommand{\cabecera}[1]{%
  \thispagestyle{plain.scrheadings}%
  \noindent
  \begin{minipage}[c]{2.2cm}
    \includegraphics[width=1.9cm]{../../comun/ucm.png}
  \end{minipage}%
  \begin{minipage}[c]{\dimexpr\linewidth-2.2cm}
    \sffamily
    {\small\color{pdsgris}\grado\par \asignatura\ · \curso\par}
    \vspace{4pt}
    {\Large\bfseries\color{pdsazul}#1\par}
    {\large\bfseries Tema~\numtema. \nombretema\par}
    \vspace{2pt}
    {\small\color{pdsgris}\autores\par}
  \end{minipage}\par
  \vspace{6pt}{\color{pdsazul}\rule{\linewidth}{0.8pt}}\par\vspace{4pt}}

% ---------------------------------------------------------------------
%  Cajas: definición, propiedad, ejemplo, nota, «en código»
% ---------------------------------------------------------------------
\tcbset{pdsbase/.style={enhanced, breakable, boxrule=0pt, arc=1.5pt, left=8pt, right=8pt,
    top=5pt, bottom=5pt, fonttitle=\sffamily\bfseries\small, coltitle=white,
    before skip=8pt, after skip=8pt}}
\newtcbtheorem[number within=section]{definicion}{Definición}%
  {pdsbase, colback=pdsfondo, colframe=pdsazul, borderline west={2.5pt}{0pt}{pdsazul}}{def}
\newtcbtheorem[use counter from=definicion]{propiedad}{Propiedad}%
  {pdsbase, colback=pdsfondo, colframe=pdsazul!75, borderline west={2.5pt}{0pt}{pdsazul!75}}{prop}
\newtcbtheorem[use counter from=definicion]{ejemplo}{Ejemplo}%
  {pdsbase, colback=pdsverde!5, colframe=pdsverde, borderline west={2.5pt}{0pt}{pdsverde}}{ej}
\newtcolorbox{nota}[1][Nota]{pdsbase, colback=pdsnaranja!7, colframe=pdsnaranja,
  title={#1}, borderline west={2.5pt}{0pt}{pdsnaranja}}
\newtcolorbox{importante}[1][Importante]{pdsbase, colback=pdsrojo!5, colframe=pdsrojo,
  title={#1}, borderline west={2.5pt}{0pt}{pdsrojo}}
\newtcolorbox{encodigo}[1][En código]{pdsbase, colback=white, colframe=pdsgris,
  title={#1}, borderline west={2.5pt}{0pt}{pdsgris}}

\lstset{basicstyle=\ttfamily\small}
\renewcommand{\codigoetiqueta}[1]{{\sffamily\small\bfseries\color{pdsgris}#1}}

% ---------------------------------------------------------------------
%  Hojas de ejercicios: \begin{ejercicio} ... \end{ejercicio} y
%  \begin{solucion} ... \end{solucion} (solo se imprime si se define
%  \SOLUCIONES antes de cargar el documento; ver *-Soluciones.tex)
% ---------------------------------------------------------------------
\newif\ifsoluciones
\ifdefined\SOLUCIONES \solucionestrue \fi
\newcounter{ejercicio}
\newenvironment{ejercicio}[1][]{%
  \refstepcounter{ejercicio}\par\vspace{6pt}%
  \noindent{\sffamily\bfseries\color{pdsazul}Ejercicio~\theejercicio}%
  \if\relax\detokenize{#1}\relax\else\ {\sffamily\color{pdsgris}(#1)}\fi\quad}%
  {\par}
\NewEnviron{solucion}{%
  \ifsoluciones
    \begin{tcolorbox}[pdsbase, colback=pdsverde!4, colframe=pdsverde,
      borderline west={2.5pt}{0pt}{pdsverde}, title={Solución}]
      \BODY
    \end{tcolorbox}
  \fi}
% \begin{incisos}[columnas] ... \end{incisos}: apartados a), b)... de un
% enunciado repartidos en columnas (3 por defecto)
% (con una sola columna no se usa multicols, que necesita al menos dos)
\newenvironment{incisos}[1][3]{\par\vspace{-2pt}%
  \ifnum#1>1 \def\finincisos{\end{multicols}}\begin{multicols}{#1}\raggedcolumns\raggedright
  \else\def\finincisos{}\fi
  \begin{enumerate}[label=\alph*), itemsep=3pt, topsep=0pt]}%
  {\end{enumerate}\finincisos\vspace{-4pt}}
\newcommand{\herramienta}[1]{\emph{Herramienta interactiva:} \texttt{herramientas/#1}}

%     \tema{1}{Introducción: señales y sistemas}
%     \begin{document} \cabecera{Apuntes} ... \end{document}
%  Compilar con:  latexmk -xelatex  (o «make» desde teoría/)
% =====================================================================

% =====================================================================
%  pds-comun.tex — Paquetes, colores y estilos de código comunes a
%  transparencias, apuntes y hojas de ejercicios. Compilar con XeLaTeX.
%  No cargar directamente: lo cargan pds-beamer.tex y pds-apuntes.tex.
% =====================================================================

\usepackage{polyglossia}
\setmainlanguage{spanish}
\usepackage{mathtools}
\usepackage{graphicx}
\usepackage{booktabs}
\usepackage{multirow}
\usepackage{array}
\usepackage{xcolor}
\usepackage{tikz}
\usetikzlibrary{positioning,arrows.meta,calc,shapes.geometric,decorations.pathreplacing}
\usepackage{listings}
\usepackage{siunitx}
\sisetup{output-decimal-marker={,}, group-separator={\,}, per-mode=symbol}

% Notas de erratas de versiones anteriores: no se imprimen salvo que se
% defina \ERRATAS antes de cargar el estilo (make ERRATAS=1)
\newif\iferratas
\ifdefined\ERRATAS \erratastrue \fi

% Ruta común (logo y figuras compartidas)
\providecommand{\rutacomun}{../../comun/}
\input{\rutacomun datos.tex}

% ---------------------------------------------------------------------
%  Paleta (la misma que usan las figuras generadas con Python,
%  ver comun/pds_estilo.py)
% ---------------------------------------------------------------------
\definecolor{pdsazul}{HTML}{2B3A8C}
\definecolor{pdsrojo}{HTML}{C0392B}
\definecolor{pdsverde}{HTML}{1E7B4F}
\definecolor{pdsnaranja}{HTML}{D9822B}
\definecolor{pdsgris}{HTML}{5A5A66}
\definecolor{pdsfondo}{HTML}{F3F4FA}
\definecolor{pdscodigo}{HTML}{F7F7F2}

% ---------------------------------------------------------------------
%  Diagramas de bloques (TikZ)
% ---------------------------------------------------------------------
\tikzset{
  bloque/.style={rectangle, rounded corners=2pt, draw=pdsazul!80, fill=pdsfondo,
                 very thick, minimum height=8mm, minimum width=14mm, align=center},
  sumador/.style={circle, draw=pdsazul!80, fill=pdsfondo, thick, inner sep=1.5pt},
  flecha/.style={-{Stealth[length=2.2mm]}, thick, pdsgris},
}

% ---------------------------------------------------------------------
%  Código: Python y Matlab en paralelo
% ---------------------------------------------------------------------
\newcommand{\pdscodesize}{\footnotesize}
\lstdefinestyle{pds}{
  basicstyle=\ttfamily\pdscodesize,
  keywordstyle=\color{pdsazul}\bfseries,
  commentstyle=\color{pdsverde}\itshape,
  stringstyle=\color{pdsrojo},
  numbers=none,
  backgroundcolor=\color{pdscodigo},
  frame=single, rulecolor=\color{pdsgris!40}, framerule=0.4pt,
  xleftmargin=2pt, xrightmargin=2pt, aboveskip=2pt, belowskip=2pt,
  columns=fullflexible, keepspaces=true,
  showstringspaces=false, breaklines=true, tabsize=4,
  upquote=true,
  % Con XeLaTeX, listings desordena los caracteres no ASCII junto a
  % dígitos («2π» sale «π2»): se sustituyen por su versión matemática.
  literate={π}{{$\pi$}}1 {δ}{{$\delta$}}1 {ω}{{$\omega$}}1 {Ω}{{$\Omega$}}1
           {≈}{{$\approx$}}1 {µ}{{$\mu$}}1 {·}{{$\cdot$}}1 {α}{{$\alpha$}}1
           {á}{{á}}1 {é}{{é}}1 {í}{{í}}1 {ó}{{ó}}1 {ú}{{ú}}1 {ñ}{{ñ}}1 {¿}{{¿}}1,
}
\lstdefinestyle{python}{style=pds, language=Python,
  morekeywords={as,True,False,None}}
\lstdefinestyle{matlab}{style=pds, language=Matlab,
  morekeywords={function,end,true,false}}

% \codigo{<ruta sin extensión>} muestra el mismo ejemplo en Python
% y en Matlab, uno al lado del otro. Los ficheros viven en
% ../codigo/python/<nombre>.py y ../codigo/matlab/<nombre>.m y se
% ejecutan en las pruebas (make prueba-codigo).
\newcommand{\rutacodigo}{../codigo/}
\newcommand{\codigoetiqueta}[1]{{\sffamily\scriptsize\bfseries\color{pdsgris}#1}}
\newcommand{\codigo}[2][]{%
  \par\noindent
  \begin{minipage}[t]{0.495\linewidth}
    \codigoetiqueta{Python}\par\vspace{1pt}
    \lstinputlisting[style=python,#1]{\rutacodigo python/#2.py}
  \end{minipage}\hfill
  \begin{minipage}[t]{0.495\linewidth}
    \codigoetiqueta{Matlab}\par\vspace{1pt}
    \lstinputlisting[style=matlab,#1]{\rutacodigo matlab/#2.m}
  \end{minipage}\par}

\input{\rutacomun notacion.tex}


\usepackage[a4paper, top=2.6cm, bottom=2.6cm, left=2.5cm, right=2.5cm,
            headheight=1.1cm]{geometry}
\usepackage{fontspec}
\usepackage{unicode-math}
\setmainfont{LibertinusSerif}[Extension=.otf, UprightFont=*-Regular, BoldFont=*-Semibold,
  ItalicFont=*-Italic, BoldItalicFont=*-SemiboldItalic]
\setsansfont{FiraSans}[Extension=.otf, UprightFont=*-Regular, BoldFont=*-SemiBold,
  ItalicFont=*-Italic, BoldItalicFont=*-SemiBoldItalic, Scale=MatchLowercase]
\setmonofont{FiraMono}[Extension=.otf, UprightFont=*-Regular, BoldFont=*-Medium, AutoFakeSlant=0.15, Scale=0.85]
\setmathfont{LibertinusMath-Regular.otf}

\usepackage{microtype}
\usepackage{enumitem}
\usepackage{multicol}
\setlist{itemsep=2pt, topsep=4pt}
\usepackage{caption}
\captionsetup{font=small, labelfont={sf,bf,color=pdsazul}, labelsep=period}
\usepackage{subcaption}
\usepackage[most]{tcolorbox}
\usepackage{environ}
\usepackage{scrlayer-scrpage}
\usepackage[hidelinks, unicode]{hyperref}
\usepackage[spanish, noabbrev]{cleveref}

% Títulos en Fira Sans con el azul de la asignatura
\setkomafont{disposition}{\sffamily\bfseries\color{pdsazul}}
\setkomafont{descriptionlabel}{\sffamily\bfseries}
\setlength{\parindent}{0pt}
\setlength{\parskip}{0.55\baselineskip plus 2pt}
\linespread{1.05}

% ---------------------------------------------------------------------
%  Datos del tema y cabeceras
% ---------------------------------------------------------------------
\newcommand{\numtema}{}
\newcommand{\nombretema}{}
\newcommand{\tema}[2]{\renewcommand{\numtema}{#1}\renewcommand{\nombretema}{#2}}

\pagestyle{scrheadings}
\clearpairofpagestyles
\ihead{\sffamily\small\color{pdsgris}\asignatura}
\ohead{\sffamily\small\color{pdsgris}Tema~\numtema. \nombretema}
\ifoot{\sffamily\small\color{pdsgris}\grado}
\ofoot{\sffamily\small\pagemark}
\setkomafont{pagehead}{\normalfont}
\setkomafont{pagefoot}{\normalfont}

% \cabecera{Apuntes|Ejercicios|...} : bloque de título de la primera página
\newcommand{\cabecera}[1]{%
  \thispagestyle{plain.scrheadings}%
  \noindent
  \begin{minipage}[c]{2.2cm}
    \includegraphics[width=1.9cm]{../../comun/ucm.png}
  \end{minipage}%
  \begin{minipage}[c]{\dimexpr\linewidth-2.2cm}
    \sffamily
    {\small\color{pdsgris}\grado\par \asignatura\ · \curso\par}
    \vspace{4pt}
    {\Large\bfseries\color{pdsazul}#1\par}
    {\large\bfseries Tema~\numtema. \nombretema\par}
    \vspace{2pt}
    {\small\color{pdsgris}\autores\par}
  \end{minipage}\par
  \vspace{6pt}{\color{pdsazul}\rule{\linewidth}{0.8pt}}\par\vspace{4pt}}

% ---------------------------------------------------------------------
%  Cajas: definición, propiedad, ejemplo, nota, «en código»
% ---------------------------------------------------------------------
\tcbset{pdsbase/.style={enhanced, breakable, boxrule=0pt, arc=1.5pt, left=8pt, right=8pt,
    top=5pt, bottom=5pt, fonttitle=\sffamily\bfseries\small, coltitle=white,
    before skip=8pt, after skip=8pt}}
\newtcbtheorem[number within=section]{definicion}{Definición}%
  {pdsbase, colback=pdsfondo, colframe=pdsazul, borderline west={2.5pt}{0pt}{pdsazul}}{def}
\newtcbtheorem[use counter from=definicion]{propiedad}{Propiedad}%
  {pdsbase, colback=pdsfondo, colframe=pdsazul!75, borderline west={2.5pt}{0pt}{pdsazul!75}}{prop}
\newtcbtheorem[use counter from=definicion]{ejemplo}{Ejemplo}%
  {pdsbase, colback=pdsverde!5, colframe=pdsverde, borderline west={2.5pt}{0pt}{pdsverde}}{ej}
\newtcolorbox{nota}[1][Nota]{pdsbase, colback=pdsnaranja!7, colframe=pdsnaranja,
  title={#1}, borderline west={2.5pt}{0pt}{pdsnaranja}}
\newtcolorbox{importante}[1][Importante]{pdsbase, colback=pdsrojo!5, colframe=pdsrojo,
  title={#1}, borderline west={2.5pt}{0pt}{pdsrojo}}
\newtcolorbox{encodigo}[1][En código]{pdsbase, colback=white, colframe=pdsgris,
  title={#1}, borderline west={2.5pt}{0pt}{pdsgris}}

\lstset{basicstyle=\ttfamily\small}
\renewcommand{\codigoetiqueta}[1]{{\sffamily\small\bfseries\color{pdsgris}#1}}

% ---------------------------------------------------------------------
%  Hojas de ejercicios: \begin{ejercicio} ... \end{ejercicio} y
%  \begin{solucion} ... \end{solucion} (solo se imprime si se define
%  \SOLUCIONES antes de cargar el documento; ver *-Soluciones.tex)
% ---------------------------------------------------------------------
\newif\ifsoluciones
\ifdefined\SOLUCIONES \solucionestrue \fi
\newcounter{ejercicio}
\newenvironment{ejercicio}[1][]{%
  \refstepcounter{ejercicio}\par\vspace{6pt}%
  \noindent{\sffamily\bfseries\color{pdsazul}Ejercicio~\theejercicio}%
  \if\relax\detokenize{#1}\relax\else\ {\sffamily\color{pdsgris}(#1)}\fi\quad}%
  {\par}
\NewEnviron{solucion}{%
  \ifsoluciones
    \begin{tcolorbox}[pdsbase, colback=pdsverde!4, colframe=pdsverde,
      borderline west={2.5pt}{0pt}{pdsverde}, title={Solución}]
      \BODY
    \end{tcolorbox}
  \fi}
% \begin{incisos}[columnas] ... \end{incisos}: apartados a), b)... de un
% enunciado repartidos en columnas (3 por defecto)
% (con una sola columna no se usa multicols, que necesita al menos dos)
\newenvironment{incisos}[1][3]{\par\vspace{-2pt}%
  \ifnum#1>1 \def\finincisos{\end{multicols}}\begin{multicols}{#1}\raggedcolumns\raggedright
  \else\def\finincisos{}\fi
  \begin{enumerate}[label=\alph*), itemsep=3pt, topsep=0pt]}%
  {\end{enumerate}\finincisos\vspace{-4pt}}
\newcommand{\herramienta}[1]{\emph{Herramienta interactiva:} \texttt{herramientas/#1}}

%     \tema{1}{Introducción: señales y sistemas}
%     \begin{document} \cabecera{Apuntes} ... \end{document}
%  Compilar con:  latexmk -xelatex  (o «make» desde teoría/)
% =====================================================================

% =====================================================================
%  pds-comun.tex — Paquetes, colores y estilos de código comunes a
%  transparencias, apuntes y hojas de ejercicios. Compilar con XeLaTeX.
%  No cargar directamente: lo cargan pds-beamer.tex y pds-apuntes.tex.
% =====================================================================

\usepackage{polyglossia}
\setmainlanguage{spanish}
\usepackage{mathtools}
\usepackage{graphicx}
\usepackage{booktabs}
\usepackage{multirow}
\usepackage{array}
\usepackage{xcolor}
\usepackage{tikz}
\usetikzlibrary{positioning,arrows.meta,calc,shapes.geometric,decorations.pathreplacing}
\usepackage{listings}
\usepackage{siunitx}
\sisetup{output-decimal-marker={,}, group-separator={\,}, per-mode=symbol}

% Notas de erratas de versiones anteriores: no se imprimen salvo que se
% defina \ERRATAS antes de cargar el estilo (make ERRATAS=1)
\newif\iferratas
\ifdefined\ERRATAS \erratastrue \fi

% Ruta común (logo y figuras compartidas)
\providecommand{\rutacomun}{../../comun/}
\input{\rutacomun datos.tex}

% ---------------------------------------------------------------------
%  Paleta (la misma que usan las figuras generadas con Python,
%  ver comun/pds_estilo.py)
% ---------------------------------------------------------------------
\definecolor{pdsazul}{HTML}{2B3A8C}
\definecolor{pdsrojo}{HTML}{C0392B}
\definecolor{pdsverde}{HTML}{1E7B4F}
\definecolor{pdsnaranja}{HTML}{D9822B}
\definecolor{pdsgris}{HTML}{5A5A66}
\definecolor{pdsfondo}{HTML}{F3F4FA}
\definecolor{pdscodigo}{HTML}{F7F7F2}

% ---------------------------------------------------------------------
%  Diagramas de bloques (TikZ)
% ---------------------------------------------------------------------
\tikzset{
  bloque/.style={rectangle, rounded corners=2pt, draw=pdsazul!80, fill=pdsfondo,
                 very thick, minimum height=8mm, minimum width=14mm, align=center},
  sumador/.style={circle, draw=pdsazul!80, fill=pdsfondo, thick, inner sep=1.5pt},
  flecha/.style={-{Stealth[length=2.2mm]}, thick, pdsgris},
}

% ---------------------------------------------------------------------
%  Código: Python y Matlab en paralelo
% ---------------------------------------------------------------------
\newcommand{\pdscodesize}{\footnotesize}
\lstdefinestyle{pds}{
  basicstyle=\ttfamily\pdscodesize,
  keywordstyle=\color{pdsazul}\bfseries,
  commentstyle=\color{pdsverde}\itshape,
  stringstyle=\color{pdsrojo},
  numbers=none,
  backgroundcolor=\color{pdscodigo},
  frame=single, rulecolor=\color{pdsgris!40}, framerule=0.4pt,
  xleftmargin=2pt, xrightmargin=2pt, aboveskip=2pt, belowskip=2pt,
  columns=fullflexible, keepspaces=true,
  showstringspaces=false, breaklines=true, tabsize=4,
  upquote=true,
  % Con XeLaTeX, listings desordena los caracteres no ASCII junto a
  % dígitos («2π» sale «π2»): se sustituyen por su versión matemática.
  literate={π}{{$\pi$}}1 {δ}{{$\delta$}}1 {ω}{{$\omega$}}1 {Ω}{{$\Omega$}}1
           {≈}{{$\approx$}}1 {µ}{{$\mu$}}1 {·}{{$\cdot$}}1 {α}{{$\alpha$}}1
           {á}{{á}}1 {é}{{é}}1 {í}{{í}}1 {ó}{{ó}}1 {ú}{{ú}}1 {ñ}{{ñ}}1 {¿}{{¿}}1,
}
\lstdefinestyle{python}{style=pds, language=Python,
  morekeywords={as,True,False,None}}
\lstdefinestyle{matlab}{style=pds, language=Matlab,
  morekeywords={function,end,true,false}}

% \codigo{<ruta sin extensión>} muestra el mismo ejemplo en Python
% y en Matlab, uno al lado del otro. Los ficheros viven en
% ../codigo/python/<nombre>.py y ../codigo/matlab/<nombre>.m y se
% ejecutan en las pruebas (make prueba-codigo).
\newcommand{\rutacodigo}{../codigo/}
\newcommand{\codigoetiqueta}[1]{{\sffamily\scriptsize\bfseries\color{pdsgris}#1}}
\newcommand{\codigo}[2][]{%
  \par\noindent
  \begin{minipage}[t]{0.495\linewidth}
    \codigoetiqueta{Python}\par\vspace{1pt}
    \lstinputlisting[style=python,#1]{\rutacodigo python/#2.py}
  \end{minipage}\hfill
  \begin{minipage}[t]{0.495\linewidth}
    \codigoetiqueta{Matlab}\par\vspace{1pt}
    \lstinputlisting[style=matlab,#1]{\rutacodigo matlab/#2.m}
  \end{minipage}\par}

\input{\rutacomun notacion.tex}


\usepackage[a4paper, top=2.6cm, bottom=2.6cm, left=2.5cm, right=2.5cm,
            headheight=1.1cm]{geometry}
\usepackage{fontspec}
\usepackage{unicode-math}
\setmainfont{LibertinusSerif}[Extension=.otf, UprightFont=*-Regular, BoldFont=*-Semibold,
  ItalicFont=*-Italic, BoldItalicFont=*-SemiboldItalic]
\setsansfont{FiraSans}[Extension=.otf, UprightFont=*-Regular, BoldFont=*-SemiBold,
  ItalicFont=*-Italic, BoldItalicFont=*-SemiBoldItalic, Scale=MatchLowercase]
\setmonofont{FiraMono}[Extension=.otf, UprightFont=*-Regular, BoldFont=*-Medium, AutoFakeSlant=0.15, Scale=0.85]
\setmathfont{LibertinusMath-Regular.otf}

\usepackage{microtype}
\usepackage{enumitem}
\usepackage{multicol}
\setlist{itemsep=2pt, topsep=4pt}
\usepackage{caption}
\captionsetup{font=small, labelfont={sf,bf,color=pdsazul}, labelsep=period}
\usepackage{subcaption}
\usepackage[most]{tcolorbox}
\usepackage{environ}
\usepackage{scrlayer-scrpage}
\usepackage[hidelinks, unicode]{hyperref}
\usepackage[spanish, noabbrev]{cleveref}

% Títulos en Fira Sans con el azul de la asignatura
\setkomafont{disposition}{\sffamily\bfseries\color{pdsazul}}
\setkomafont{descriptionlabel}{\sffamily\bfseries}
\setlength{\parindent}{0pt}
\setlength{\parskip}{0.55\baselineskip plus 2pt}
\linespread{1.05}

% ---------------------------------------------------------------------
%  Datos del tema y cabeceras
% ---------------------------------------------------------------------
\newcommand{\numtema}{}
\newcommand{\nombretema}{}
\newcommand{\tema}[2]{\renewcommand{\numtema}{#1}\renewcommand{\nombretema}{#2}}

\pagestyle{scrheadings}
\clearpairofpagestyles
\ihead{\sffamily\small\color{pdsgris}\asignatura}
\ohead{\sffamily\small\color{pdsgris}Tema~\numtema. \nombretema}
\ifoot{\sffamily\small\color{pdsgris}\grado}
\ofoot{\sffamily\small\pagemark}
\setkomafont{pagehead}{\normalfont}
\setkomafont{pagefoot}{\normalfont}

% \cabecera{Apuntes|Ejercicios|...} : bloque de título de la primera página
\newcommand{\cabecera}[1]{%
  \thispagestyle{plain.scrheadings}%
  \noindent
  \begin{minipage}[c]{2.2cm}
    \includegraphics[width=1.9cm]{../../comun/ucm.png}
  \end{minipage}%
  \begin{minipage}[c]{\dimexpr\linewidth-2.2cm}
    \sffamily
    {\small\color{pdsgris}\grado\par \asignatura\ · \curso\par}
    \vspace{4pt}
    {\Large\bfseries\color{pdsazul}#1\par}
    {\large\bfseries Tema~\numtema. \nombretema\par}
    \vspace{2pt}
    {\small\color{pdsgris}\autores\par}
  \end{minipage}\par
  \vspace{6pt}{\color{pdsazul}\rule{\linewidth}{0.8pt}}\par\vspace{4pt}}

% ---------------------------------------------------------------------
%  Cajas: definición, propiedad, ejemplo, nota, «en código»
% ---------------------------------------------------------------------
\tcbset{pdsbase/.style={enhanced, breakable, boxrule=0pt, arc=1.5pt, left=8pt, right=8pt,
    top=5pt, bottom=5pt, fonttitle=\sffamily\bfseries\small, coltitle=white,
    before skip=8pt, after skip=8pt}}
\newtcbtheorem[number within=section]{definicion}{Definición}%
  {pdsbase, colback=pdsfondo, colframe=pdsazul, borderline west={2.5pt}{0pt}{pdsazul}}{def}
\newtcbtheorem[use counter from=definicion]{propiedad}{Propiedad}%
  {pdsbase, colback=pdsfondo, colframe=pdsazul!75, borderline west={2.5pt}{0pt}{pdsazul!75}}{prop}
\newtcbtheorem[use counter from=definicion]{ejemplo}{Ejemplo}%
  {pdsbase, colback=pdsverde!5, colframe=pdsverde, borderline west={2.5pt}{0pt}{pdsverde}}{ej}
\newtcolorbox{nota}[1][Nota]{pdsbase, colback=pdsnaranja!7, colframe=pdsnaranja,
  title={#1}, borderline west={2.5pt}{0pt}{pdsnaranja}}
\newtcolorbox{importante}[1][Importante]{pdsbase, colback=pdsrojo!5, colframe=pdsrojo,
  title={#1}, borderline west={2.5pt}{0pt}{pdsrojo}}
\newtcolorbox{encodigo}[1][En código]{pdsbase, colback=white, colframe=pdsgris,
  title={#1}, borderline west={2.5pt}{0pt}{pdsgris}}

\lstset{basicstyle=\ttfamily\small}
\renewcommand{\codigoetiqueta}[1]{{\sffamily\small\bfseries\color{pdsgris}#1}}

% ---------------------------------------------------------------------
%  Hojas de ejercicios: \begin{ejercicio} ... \end{ejercicio} y
%  \begin{solucion} ... \end{solucion} (solo se imprime si se define
%  \SOLUCIONES antes de cargar el documento; ver *-Soluciones.tex)
% ---------------------------------------------------------------------
\newif\ifsoluciones
\ifdefined\SOLUCIONES \solucionestrue \fi
\newcounter{ejercicio}
\newenvironment{ejercicio}[1][]{%
  \refstepcounter{ejercicio}\par\vspace{6pt}%
  \noindent{\sffamily\bfseries\color{pdsazul}Ejercicio~\theejercicio}%
  \if\relax\detokenize{#1}\relax\else\ {\sffamily\color{pdsgris}(#1)}\fi\quad}%
  {\par}
\NewEnviron{solucion}{%
  \ifsoluciones
    \begin{tcolorbox}[pdsbase, colback=pdsverde!4, colframe=pdsverde,
      borderline west={2.5pt}{0pt}{pdsverde}, title={Solución}]
      \BODY
    \end{tcolorbox}
  \fi}
% \begin{incisos}[columnas] ... \end{incisos}: apartados a), b)... de un
% enunciado repartidos en columnas (3 por defecto)
% (con una sola columna no se usa multicols, que necesita al menos dos)
\newenvironment{incisos}[1][3]{\par\vspace{-2pt}%
  \ifnum#1>1 \def\finincisos{\end{multicols}}\begin{multicols}{#1}\raggedcolumns\raggedright
  \else\def\finincisos{}\fi
  \begin{enumerate}[label=\alph*), itemsep=3pt, topsep=0pt]}%
  {\end{enumerate}\finincisos\vspace{-4pt}}
\newcommand{\herramienta}[1]{\emph{Herramienta interactiva:} \texttt{herramientas/#1}}

\tema{2}{Muestreo y reconstrucción de señales}

\begin{document}
\ifsoluciones\cabecera{Ejercicios resueltos}\else\cabecera{Ejercicios}\fi

Los ejercicios 1–6 corresponden a la parte~1 del tema (muestreo y aliasing), los ejercicios 7–10 a la parte~2 (reconstrucción), los ejercicios 11–13 a la parte~3 (cuantificación y codificación) y los ejercicios 14–16 a la parte~4 (muestreo de señales discretas). Se usa $\sinc(x)=\sin(\pi x)/(\pi x)$ y $\Sa(x)=\sin x/x$. Muchos resultados se pueden comprobar numéricamente con los ejemplos de \texttt{teoría/T2/codigo/}.

% =====================================================================
\section*{Muestreo y aliasing}

\begin{ejercicio}[frecuencias discretas]
  La señal $x(t)=3\cos(50\pi t)+10\sin(300\pi t)-\cos(100\pi t)$ ($t$ en segundos) se muestrea con $f_s=200$\,Hz.
  \begin{incisos}[1]
    \item ¿Cuál es la tasa de Nyquist de $x(t)$?
    \item Obtén la secuencia $x[n]$ expresando cada término con su alias fundamental, $\abs{\Omega}\le\pi$.
    \item ¿Qué señal se obtiene si se reconstruye $x[n]$ con el filtro ideal de corte $f_s/2$?
  \end{incisos}
\end{ejercicio}
\begin{solucion}
  \begin{enumerate}[label=\alph*)]
    \item Las frecuencias son 25, 150 y 50\,Hz. La máxima es $f_m=150$\,Hz, así que la tasa de Nyquist es $2f_m=300$\,Hz.
    \item $\Omega=2\pi f/f_s$ da $\pi/4$, $3\pi/2$ y $\pi/2$. La segunda se pliega a $3\pi/2-2\pi=-\pi/2$, y como el seno es impar,
    \[
      x[n]=3\cos\big(\tfrac\pi4n\big)-10\sin\big(\tfrac\pi2n\big)-\cos\big(\tfrac\pi2n\big).
    \]
    \item Cada $\Omega$ vuelve a $f=\Omega f_s/2\pi$: $y(t)=3\cos(50\pi t)-10\sin(100\pi t)-\cos(100\pi t)$. El tono de 150\,Hz aparece en 50\,Hz con el signo cambiado y se suma con el de 50\,Hz original.
  \end{enumerate}
\end{solucion}

\begin{ejercicio}[alias fundamental]
  Las sinusoides $x(t)=\sin(2\pi f_0t)$ con $f_0=200$, $800$, $1200$ y $1700$\,Hz se muestrean a $f_s=1$\,kHz. Para cada una, calcula la frecuencia discreta de su alias fundamental y la señal que se obtendría al reconstruir con el filtro ideal de corte $f_s/2$.
\end{ejercicio}
\begin{solucion}
  Con $\Omega_0=2\pi f_0/f_s$ llevada a $[-\pi,\pi)$:
  \begin{center}
    \begin{tabular}{@{}rlll@{}}
      \toprule
      $f_0$ (Hz) & $\Omega_0$ & alias fundamental & señal reconstruida\\\midrule
      200 & $0{,}4\pi$ & $0{,}4\pi$ & $\sin(400\pi t)$\\
      800 & $1{,}6\pi$ & $-0{,}4\pi$ & $-\sin(400\pi t)$\\
      1200 & $2{,}4\pi$ & $0{,}4\pi$ & $\sin(400\pi t)$\\
      1700 & $3{,}4\pi$ & $-0{,}6\pi$ & $-\sin(600\pi t)$\\\bottomrule
    \end{tabular}
  \end{center}
  Las de 200, 800 y 1200\,Hz dan las mismas muestras salvo el signo; la de 1700\,Hz aparece en 300\,Hz. Un alias negativo invierte el signo de un seno (no el de un coseno). Se comprueba con el ejemplo \texttt{alias\_fundamental}.
\end{solucion}

\begin{ejercicio}[ancho de banda de señales derivadas]
  La señal $x(t)$ tiene banda limitada a $\omega_m$. Determina la frecuencia de muestreo mínima para muestrear sin aliasing:
  \begin{incisos}[2]
    \item $x(t)+x(t-1)$
    \item $\dfrac{\dd x(t)}{\dd t}$
    \item $x^2(t)$
    \item $x(t)\cos(\omega_0t)$, con $\omega_0>\omega_m$
  \end{incisos}
\end{ejercicio}
\begin{solucion}
  \begin{enumerate}[label=\alph*)]
    \item $X(j\omega)(1+\e^{-j\omega})$: multiplicar por una función no añade frecuencias. La banda sigue siendo $\omega_m$ (como mucho, se anulan algunas): $\omega_s>2\omega_m$.
    \item $j\omega X(j\omega)$: misma banda, $\omega_s>2\omega_m$.
    \item El producto en el tiempo es una convolución en frecuencia, $\frac{1}{2\pi}X\conv X$, que ocupa $\abs{\omega}\le2\omega_m$: $\omega_s>4\omega_m$.
    \item $\frac12\big[X(j(\omega-\omega_0))+X(j(\omega+\omega_0))\big]$: el espectro se traslada a $\pm\omega_0$ y ocupa $\omega_0-\omega_m\le\abs{\omega}\le\omega_0+\omega_m$. La frecuencia máxima es $\omega_0+\omega_m$, así que el teorema del muestreo exige $\omega_s>2(\omega_0+\omega_m)$. \emph{No} basta con $2\omega_m$: con esa frecuencia las réplicas de las bandas en $\pm\omega_0$ se solapan. (Para señales paso banda existe un muestreo paso banda que permite frecuencias menores con condiciones adicionales; queda fuera de este tema.)
  \end{enumerate}
\end{solucion}

\begin{ejercicio}[periodo de muestreo]
  Determina el periodo de muestreo necesario para poder reconstruir las siguientes señales. Si no es posible evitar el aliasing, propón un criterio práctico.
  \begin{incisos}[3]
    \item $x_1(t)=\e^{-t}u(t)$
    \item $x_2(t)=\sinc(f_ct)$
    \item $x_3(t)=\sinc^2(f_ct)$
  \end{incisos}
\end{ejercicio}
\begin{solucion}
  \begin{enumerate}[label=\alph*)]
    \item $X_1(j\omega)=1/(1+j\omega)$ no se anula nunca: la señal no es de banda limitada y siempre hay aliasing. Como $\abs{X_1}=1/\sqrt{1+\omega^2}$ decae, se puede elegir $\omega_s/2$ donde $\abs{X_1}$ sea despreciable. Por ejemplo, para que $\abs{X_1}<0{,}01\,\abs{X_1(0)}$ hace falta $\omega>\sqrt{10^4-1}\approx100$\,rad/s, así que $\omega_s\approx200$\,rad/s y $T_s\approx0{,}03$\,s.
    \item $\sinc(f_ct)\parF\frac{1}{f_c}\rect(f/f_c)$, que ocupa $\abs{f}\le f_c/2$. La banda es $f_m=f_c/2$, no $f_c$: hace falta $f_s>f_c$, es decir, $T_s<1/f_c$.
    \item Su transformada es la convolución de dos pulsos, un triángulo que ocupa $\abs{f}\le f_c$: $f_m=f_c$, $f_s>2f_c$ y $T_s<1/(2f_c)$. El ejemplo \texttt{espectro\_muestreado} usa esta señal con $f_c=1$\,Hz y $f_s=3$\,Hz.
  \end{enumerate}
\end{solucion}

\begin{ejercicio}[espectro de la señal muestreada]
  Una señal $x(t)$ tiene espectro triangular: $X(f)=1-\abs{f}/100$ para $\abs{f}\le100$\,Hz y cero fuera.
  \begin{incisos}[1]
    \item Dibuja el espectro de $x_s(t)$ para $f_s=300$\,Hz y para $f_s=150$\,Hz.
    \item Con $f_s=150$\,Hz, ¿en qué banda hay aliasing? ¿Qué forma tiene allí el espectro? ¿Qué banda se conserva intacta?
    \item Si antes de muestrear a 150\,Hz se aplica un filtro antialiasing ideal, ¿cuál debe ser su corte y qué se pierde?
  \end{incisos}
\end{ejercicio}
\begin{solucion}
  \begin{enumerate}[label=\alph*)]
    \item $X_s(f)=f_s\sum_kX(f-kf_s)$: triángulos de altura $f_s$ centrados en $kf_s$. Con 300\,Hz, cada triángulo ocupa $[k\cdot300-100,k\cdot300+100]$ y quedan huecos de 100\,Hz entre ellos. Con 150\,Hz, cada triángulo se solapa con sus vecinos.
    \item La réplica en 150\,Hz ocupa $[50,250]$ y se solapa con la central en $50\le f\le100$. Allí la suma vale $f_s\big[(1-\frac{f}{100})+(1-\frac{150-f}{100})\big]=0{,}5\,f_s$, constante: el espectro se aplana. La banda $\abs{f}<50$\,Hz se conserva intacta.
    \item Corte en $f_s/2=75$\,Hz. Entonces $X$ ocupa $\abs{f}\le75$, las réplicas empiezan en $150-75=75$\,Hz y no se solapan. La banda $\abs{f}<75$\,Hz se conserva y se pierde la de 75 a 100\,Hz: es preferible perderla a que se mezcle con la banda útil.
  \end{enumerate}
\end{solucion}

\begin{ejercicio}[frecuencias en un electrocardiograma]
  Un electrocardiograma (ECG), con información útil entre 0,05 y 40\,Hz, se muestrea a $f_s=500$\,Hz. La red eléctrica añade una interferencia de 50\,Hz y sus armónicos.
  \begin{incisos}[1]
    \item ¿A qué frecuencias discretas $\Omega$ aparecen la interferencia y sus armónicos de 150 y 250\,Hz?
    \item Un filtro digital tiene un cero en $\Omega=0{,}1\pi$. ¿Qué frecuencia en hercios elimina?
    \item Si el mismo ECG se muestreara a 120\,Hz sin filtro antialiasing, ¿dónde aparecería el armónico de 150\,Hz? ¿Se podría eliminar después con un filtro digital?
  \end{incisos}
\end{ejercicio}
\begin{solucion}
  \begin{enumerate}[label=\alph*)]
    \item $\Omega=2\pi f/f_s$: 50\,Hz en $\pi/5$, 150\,Hz en $3\pi/5$ y 250\,Hz en $\pi$, que es la frecuencia de Nyquist.
    \item $f=\Omega f_s/2\pi=0{,}05\cdot500=25$\,Hz.
    \item $150/120=1{,}25$, así que el entero más próximo es $k=1$ y $f_a=\abs{150-120}=30$\,Hz, es decir, $\Omega=\pi/2$. Cae dentro de la banda útil del ECG: un filtro digital que lo eliminara se llevaría también la parte de la señal en 30\,Hz. Por eso el filtro antialiasing es imprescindible antes de muestrear.
  \end{enumerate}
\end{solucion}

% =====================================================================
\section*{Reconstrucción}

\begin{ejercicio}[reconstrucción ideal]
  La señal $x(t)=\cos(6\pi t)+\cos(14\pi t)$ se muestrea y se reconstruye con un filtro paso bajo ideal.
  \begin{incisos}[1]
    \item Con $f_s=12$\,Hz y corte $f_s/2$, ¿qué señal se obtiene?
    \item ¿Cuál es la frecuencia de muestreo mínima para reconstruirla exactamente?
    \item Con $f_s=20$\,Hz, ¿entre qué valores puede estar la frecuencia de corte del filtro?
  \end{incisos}
\end{ejercicio}
\begin{solucion}
  \begin{enumerate}[label=\alph*)]
    \item Las frecuencias son 3 y 7\,Hz. La de 7\,Hz supera $f_s/2=6$\,Hz y se pliega a $\abs{7-12}=5$\,Hz; al ser un coseno no cambia de signo. Se obtiene $y(t)=\cos(6\pi t)+\cos(10\pi t)$, que tiene las mismas muestras que $x(t)$ a 12\,Hz.
    \item $f_s>2\cdot7=14$\,Hz.
    \item Hay que dejar pasar 7\,Hz y eliminar la réplica más baja, en $20-7=13$\,Hz: $7<f_c<13$\,Hz.
  \end{enumerate}
\end{solucion}

\begin{ejercicio}[retenedores de orden uno]
  Obtén la respuesta al impulso y la respuesta en frecuencia de los tres retenedores de orden uno:
  \begin{incisos}[1]
    \item El no causal, que une las muestras $x[n]$ y $x[n+1]$ con una recta.
    \item El mismo retrasado una muestra, para hacerlo causal.
    \item El causal predictivo, que prolonga la recta que pasa por $x[n-1]$ y $x[n]$.
  \end{incisos}
  Compara sus espectros con el del retenedor de orden cero. \emph{Sugerencia}: escribe $h(t)$ como convolución o suma de pulsos rectangulares.
\end{ejercicio}
\begin{solucion}
  \begin{enumerate}[label=\alph*)]
    \item Una muestra aislada $x[0]=1$ produce un triángulo que vale 1 en $t=0$ y 0 en $t=\pm T_s$: $h_1(t)=\tri(t/T_s)$. Como $\tri(t/T_s)=\frac1{T_s}\rect(t/T_s)\conv\rect(t/T_s)$ y $\rect(t/T_s)\parF T_s\sinc(\omega/\omega_s)$,
    \[
      H_1(j\omega)=\frac{1}{T_s}\big(T_s\sinc(\omega/\omega_s)\big)^2=T_s\,\sinc^2\Big(\frac{\omega}{\omega_s}\Big).
    \]
    Es real y no negativa (fase nula).
    \item $h(t)=\tri\big(\frac{t-T_s}{T_s}\big)$, no nula en $[0,2T_s]$, y $H(j\omega)=T_s\,\e^{-j\omega T_s}\sinc^2(\omega/\omega_s)$: el mismo módulo con un retardo de una muestra.
    \item $y(t)=x[n]+(x[n]-x[n-1])\frac{t-nT_s}{T_s}$ en $[nT_s,(n+1)T_s)$. La muestra $x[0]=1$ interviene como $x[n]$ en $[0,T_s)$, con contribución $1+t/T_s$, y como $x[n-1]$ en $[T_s,2T_s)$, con $-(t-T_s)/T_s=1-t/T_s$:
    \[
      h(t)=\begin{cases}1+t/T_s & 0\le t<T_s\\1-t/T_s & T_s\le t<2T_s\\0&\text{en otro caso}\end{cases}
    \]
    Integrando, $H(j\omega)=T_s(1+j\omega T_s)\,\e^{-j\omega T_s}\sinc^2(\omega/\omega_s)$ (comprobado simbólicamente).
  \end{enumerate}
  El ZOH tiene $\abs{H_0}=T_s\abs{\sinc(\omega/\omega_s)}$, con réplicas que caen como $1/\omega$; los FOH caen como $1/\omega^2$ y las atenúan mejor. El predictivo tiene un factor $\abs{1+j\omega T_s}$ que realza las altas frecuencias: su salida tiene saltos.
  \begin{nota}[Errata de versiones anteriores]
    Una versión anterior de esta hoja daba para el FOH causal un pulso en rampa de anchura $T_s$, que no reproduce la interpolación lineal, y $H=\e^{-j\omega T_s/2}\sinc^2$. Las expresiones correctas son las de arriba.
  \end{nota}
\end{solucion}

\begin{ejercicio}[conversor D/A de audio]
  Un conversor D/A de audio trabaja a $f_s=44{,}1$\,kHz con un retenedor de orden cero.
  \begin{incisos}[1]
    \item ¿Cuánto atenúa el ZOH los tonos de 1, 10 y 20\,kHz?
    \item Para un tono de 20\,kHz, ¿a qué frecuencia aparece la primera imagen y a qué nivel respecto al tono?
    \item ¿Qué ganancia debe tener un filtro digital de compensación en 20\,kHz?
    \item Repite b) si antes del conversor se interpola por 4, de modo que el ZOH trabaja a 176,4\,kHz.
  \end{incisos}
\end{ejercicio}
\begin{solucion}
  $\abs{H_0}/T_s=\abs{\sinc(f/f_s)}$.
  \begin{enumerate}[label=\alph*)]
    \item $-0{,}007$\,dB en 1\,kHz, $-0{,}75$\,dB en 10\,kHz y $-3{,}17$\,dB en 20\,kHz.
    \item La réplica en $f_s$ de la componente en $-20$\,kHz cae en $44{,}1-20=24{,}1$\,kHz, atenuada $-4{,}79$\,dB: solo $1{,}6$\,dB por debajo del tono. El ZOH no basta; hace falta un filtro de suavizado que pase 20\,kHz y elimine 24,1\,kHz, muy abrupto.
    \item $+3{,}17$\,dB, es decir, $1/\sinc(20/44{,}1)\approx1{,}44$.
    \item Ahora el tono se atenúa solo $0{,}18$\,dB y la primera imagen está en $176{,}4-20=156{,}4$\,kHz, $17{,}9$\,dB por debajo del tono. El filtro analógico de suavizado tiene una transición de 20 a 156\,kHz y puede ser muy sencillo.
  \end{enumerate}
\end{solucion}

\begin{ejercicio}[reconstrucción con retenedores]
  Muestrea las siguientes señales, eligiendo un valor adecuado de $T_s$. Indica en cada caso si es posible evitar el aliasing y estima el error de reconstruir con retenedores de orden cero y uno.
  \begin{incisos}[1]
    \item $x_1(t)=\sin\big(\frac{2\pi}{4}t\big)$
    \item $x_2(t)=\sin\big(\frac{2\pi}{4}t\big)$ para $0\le t\le7$ y cero fuera
    \item $x_3(t)=\sin\big(\frac{20\pi}{700}t\big)\cos\big(\frac{20\pi}{35}t\big)$
  \end{incisos}
  \emph{Sugerencia}: para una señal suave, el error del ZOH es como mucho $\max\abs{x'}\,T_s$ y el de la interpolación lineal, $\max\abs{x''}\,T_s^2/8$. Comprueba los resultados modificando el ejemplo \texttt{retenedores}.
\end{ejercicio}
\begin{solucion}
  \begin{enumerate}[label=\alph*)]
    \item $f_0=0{,}25$\,Hz: basta $T_s<2$\,s, aunque con $T_s$ cerca de ese límite los retenedores dan una reconstrucción muy pobre. Con $T_s=0{,}25$\,s (16 muestras por periodo) y $\omega_0=\pi/2$, el error del ZOH es como mucho $\omega_0T_s\approx0{,}39$ y el del FOH, $\omega_0^2T_s^2/8\approx0{,}019$; la simulación da $0{,}38$ y $0{,}019$. El ZOH necesita muchas más muestras para el mismo error.
    \item Es una señal de duración finita, así que no tiene banda limitada: el aliasing no se puede evitar del todo. Su espectro está concentrado alrededor de $0{,}25$\,Hz y decae, de modo que con el mismo $T_s=0{,}25$\,s el aliasing es pequeño. Los errores son los del apartado anterior, salvo cerca de $t=7$, donde la señal salta de $\sin(3{,}5\pi)=-1$ a 0 y cualquier retenedor comete un error del orden del salto.
    \item Por la fórmula del producto, $x_3(t)=\frac12\big[\sin\big(\frac{3\pi}{5}t\big)-\sin\big(\frac{19\pi}{35}t\big)\big]$, con frecuencias $0{,}3$ y $19/70\approx0{,}27$\,Hz. Hace falta $f_s>0{,}6$\,Hz, $T_s<5/3$\,s. La envolvente es lenta (0,014\,Hz) y la portadora, de $2/7$\,Hz; con $T_s=0{,}2$\,s, unas 17 muestras por periodo de portadora, el FOH reproduce bien la señal.
  \end{enumerate}
\end{solucion}

% =====================================================================
\section*{Cuantificación y codificación}

\begin{ejercicio}[cuantificador de 3 bits]
  Un cuantificador de $b=3$ bits tiene margen $[-1,1)$.
  \begin{incisos}[1]
    \item Calcula $\Delta$ y los niveles.
    \item Cuantifica $x[n]=0{,}95\cos(\pi n/4)$ para $n=0,\dots,7$ por redondeo y por truncamiento, y calcula el error.
    \item Codifica en complemento a dos el resultado del redondeo.
  \end{incisos}
\end{ejercicio}
\begin{solucion}
  \begin{enumerate}[label=\alph*)]
    \item $\Delta=2/2^3=0{,}25$; los niveles son $k\Delta$ con $k=-4,\dots,3$, de $-1$ a $0{,}75$.
    \item Los valores son $0{,}95$, $0{,}672$, $0$, $-0{,}672$, $-0{,}95$, $-0{,}672$, $0$, $0{,}672$. Para $n=0,\dots,4$ (el resto se repite por simetría):
    \begin{center}
      \begin{tabular}{@{}rrrrrr@{}}
        \toprule
        $n$ & $x[n]$ & redondeo & $e$ & truncamiento & $e$\\\midrule
        0 & $0{,}950$ & $0{,}75$ & $-0{,}200$ & $0{,}75$ & $-0{,}200$\\
        1 & $0{,}672$ & $0{,}75$ & $0{,}078$ & $0{,}50$ & $-0{,}172$\\
        2 & $0$ & $0$ & $0$ & $0$ & $0$\\
        3 & $-0{,}672$ & $-0{,}75$ & $-0{,}078$ & $-0{,}75$ & $-0{,}078$\\
        4 & $-0{,}950$ & $-1{,}00$ & $-0{,}050$ & $-1{,}00$ & $-0{,}050$\\\bottomrule
      \end{tabular}
    \end{center}
    En $n=0$ el redondeo daría $1{,}0$, que no es un nivel: el cuantificador \emph{satura} en $0{,}75$ y el error, $-0{,}2$, supera $\Delta/2=0{,}125$. El margen bipolar es asimétrico: llega a $-1$ pero no a $+1$.
    \item $k=3,3,0,-3,-4,-3,0,3$: \texttt{011}, \texttt{011}, \texttt{000}, \texttt{101}, \texttt{100}, \texttt{101}, \texttt{000}, \texttt{011}.
  \end{enumerate}
\end{solucion}

\begin{ejercicio}[conversor A/D de un microcontrolador]
  Un conversor A/D de 10 bits admite una entrada entre 0 y $V_{\text{ref}}$, que se puede configurar en 1,1, 3,3 o 5\,V.
  \begin{incisos}[1]
    \item Indica la resolución, el número de niveles y el rango dinámico de cada configuración.
    \item ¿Qué configuración es la más adecuada si la entrada varía entre 0 y 3\,V? ¿Y entre 0 y 5\,V? ¿Y entre $-2{,}5$ y $2{,}5$\,V?
    \item ¿Cuántos bits harían falta, en cada configuración, para una resolución de $10^{-4}$\,V?
  \end{incisos}
\end{ejercicio}
\begin{solucion}
  \begin{enumerate}[label=\alph*)]
    \item Siempre hay $2^{10}=1024$ niveles y $\mathrm{DR}=20\log_{10}1024\approx60{,}2$\,dB. La resolución es $\Delta=V_{\text{ref}}/2^{10}$: $1{,}07$\,mV, $3{,}22$\,mV y $4{,}88$\,mV.
    \item De 0 a 3\,V, $V_{\text{ref}}=3{,}3$\,V: cubre la señal sin saturar y con la menor $\Delta$ posible (con 1,1\,V saturaría y con 5\,V se desperdiciaría el 40\,\% de los niveles). De 0 a 5\,V, $V_{\text{ref}}=5$\,V. Entre $-2{,}5$ y $2{,}5$\,V ninguna sirve directamente: hay que sumar 2,5\,V con un circuito de adaptación (o atenuar y desplazar) y usar $V_{\text{ref}}=5$\,V, o un conversor bipolar.
    \item $b\ge\log_2(V_{\text{ref}}/10^{-4})$: $13{,}43$, $15{,}01$ y $15{,}61$, es decir, 14, 16 y 16 bits.
  \end{enumerate}
\end{solucion}

\begin{ejercicio}[relación señal a ruido de cuantificación]
  \begin{incisos}[1]
    \item ¿Cuántos bits hacen falta para que una sinusoide que ocupa todo el margen tenga $\mathrm{SQNR}\ge80$\,dB?
    \item ¿Y si su amplitud es solo la cuarta parte del margen?
    \item ¿Qué SQNR tiene una grabación de 16 bits a fondo de escala?
    \item Calcula la SQNR de una señal cuyas muestras se distribuyen uniformemente en todo el margen $[-V,V)$.
  \end{incisos}
\end{ejercicio}
\begin{solucion}
  \begin{enumerate}[label=\alph*)]
    \item $6{,}02\,b+1{,}76\ge80$ da $b\ge12{,}997$: 13 bits.
    \item La potencia baja en un factor 16, es decir, $12{,}04$\,dB, que hay que compensar: $b\ge(80+12{,}04-1{,}76)/6{,}02=14{,}997$, 15 bits.
    \item $6{,}02\cdot16+1{,}76\approx98{,}1$\,dB.
    \item La potencia de una variable uniforme en $[-V,V)$ es $P_x=(2V)^2/12$, y $\sigma_e^2=\Delta^2/12=(2V)^2/(12\cdot2^{2b})$. Entonces $\mathrm{SQNR}=10\log_{10}2^{2b}\approx6{,}02\,b$\,dB: $1{,}76$\,dB menos que la sinusoide, porque su factor de cresta es mayor.
  \end{enumerate}
\end{solucion}

% =====================================================================
\section*{Muestreo de señales discretas}

\begin{ejercicio}[diezmado]
  Sea $x[n]=\cos(0{,}2\pi n)+\cos(0{,}3\pi n)$.
  \begin{incisos}[1]
    \item ¿Cuál es el máximo factor de diezmado sin aliasing?
    \item Obtén $x[2n]$ y $x[4n]$. ¿Hay aliasing en algún caso?
    \item Antes de diezmar por 4 se aplica un paso bajo ideal de corte $\pi/4$. ¿Qué secuencia se obtiene?
  \end{incisos}
\end{ejercicio}
\begin{solucion}
  \begin{enumerate}[label=\alph*)]
    \item $\Omega_m=0{,}3\pi<\pi/M$ exige $M<3{,}33$: $M=3$, que da $\cos(0{,}6\pi n)+\cos(0{,}9\pi n)$.
    \item $x[2n]=\cos(0{,}4\pi n)+\cos(0{,}6\pi n)$, sin aliasing. $x[4n]=\cos(0{,}8\pi n)+\cos(1{,}2\pi n)$, y $1{,}2\pi$ se pliega a $-0{,}8\pi$: $x[4n]=2\cos(0{,}8\pi n)$. Las dos componentes se han mezclado y ya no se pueden separar.
    \item El filtro elimina $\cos(0{,}3\pi n)$, porque $0{,}3\pi>\pi/4$, y deja $\cos(0{,}2\pi n)$. Después del diezmado queda $\cos(0{,}8\pi n)$: se pierde una componente, pero la otra no se contamina.
  \end{enumerate}
\end{solucion}

\begin{ejercicio}[expansión e interpolación]
  \begin{incisos}[1]
    \item Se expande $x[n]=\cos(0{,}4\pi n)$ por $L=2$. ¿Qué frecuencias contiene $x_e[n]$ en $[0,\pi]$? ¿Qué se obtiene tras un filtro ideal de ganancia 2 y corte $\pi/2$?
    \item Una secuencia tiene espectro triangular de banda $\Omega_m=\pi/2$. Dibuja el espectro tras expandirla por $L=3$ e indica dónde están las imágenes.
  \end{incisos}
\end{ejercicio}
\begin{solucion}
  \begin{enumerate}[label=\alph*)]
    \item $x_e[n]=\cos(0{,}2\pi n)\cdot\frac{1+(-1)^n}{2}$, ya que vale $\cos(0{,}4\pi\,n/2)$ en los pares y 0 en los impares. Como $(-1)^n\cos(0{,}2\pi n)=\cos(1{,}2\pi n)=\cos(0{,}8\pi n)$,
    \[
      x_e[n]=\tfrac12\cos(0{,}2\pi n)+\tfrac12\cos(0{,}8\pi n).
    \]
    Es lo que predice $X(\e^{j2\Omega})$: picos donde $2\Omega=\pm0{,}4\pi+2\pi k$, es decir, en $0{,}2\pi$ y en la imagen $0{,}8\pi$. El filtro elimina la imagen y duplica el resto: $y[n]=\cos(0{,}2\pi n)$, que es $x[n]$ a doble frecuencia de muestreo.
    \item $X(\e^{j3\Omega})$ repite el triángulo cada $2\pi/3$ con banda $\pi/6$: el original queda en $\abs{\Omega}\le\pi/6$ y las imágenes, centradas en $\pm2\pi/3$, ocupan $\pi/2\le\abs{\Omega}\le5\pi/6$. El filtro de interpolación, con corte $\pi/3$ y ganancia 3, deja solo la banda central.
  \end{enumerate}
\end{solucion}

\begin{ejercicio}[cambio racional y sobremuestreo]
  \begin{incisos}[1]
    \item Diseña el esquema para pasar audio de 48\,kHz a 32\,kHz: factores $L$ y $M$, frecuencia intermedia y frecuencia de corte del filtro, en radianes por muestra y en hercios.
    \item Un conversor A/D de 12 bits sobremuestrea por $K=16$ y después diezma. ¿Qué SQNR alcanza para una sinusoide a fondo de escala? ¿A cuántos bits equivale?
    \item ¿Qué factor de sobremuestreo haría falta para ganar 4 bits?
  \end{incisos}
\end{ejercicio}
\begin{solucion}
  \begin{enumerate}[label=\alph*)]
    \item $32/48=2/3$: $L=2$, $M=3$. Se interpola a 96\,kHz y se filtra con ganancia 2 y corte $\min(\pi/2,\pi/3)=\pi/3$, que a 96\,kHz son $\frac{1}{6}\cdot96=16$\,kHz, la frecuencia de Nyquist de la salida. Después se diezma por 3.
    \item $6{,}02\cdot12+1{,}76+10\log_{10}16=72{,}0+12{,}0\approx86{,}0$\,dB, unos 2 bits más: equivale a un conversor de 14 bits.
    \item 4 bits son unos 24\,dB, así que $10\log_{10}K=24$ y $K\approx256$. Ganar resolución solo con sobremuestreo es caro; los conversores sigma-delta lo consiguen con mucho menos gracias al conformado de ruido.
  \end{enumerate}
\end{solucion}

\end{document}
